\documentclass[10pt]{article}
\usepackage[margin=0.8in]{geometry}
\usepackage[T1]{fontenc}
\usepackage{lmodern}
\usepackage{xcolor,soul}\sethlcolor{yellow}
\usepackage[normalem]{ulem}
\setlength{\parindent}{0pt}
\setlength{\parskip}{0.35em}
\pagestyle{empty}
% the two facts only the author knows, highlighted until filled in
\newcommand{\tofill}[1]{\hl{#1}}
\newcommand{\new}[1]{{\color{red}#1}}   % what changed since the author's comments
\newcommand{\del}[1]{{\color{red}\sout{#1}}}   % what was removed
\begin{document}

Professor Lawrence F.\ Katz\\
Editor, \textit{The Quarterly Journal of Economics}

Dear Professor Katz,

When you declined my paper ``Housing First or Treatment First?'' (MS 45355), you did not see a large enough advance or a sufficiently compelling empirical base. I have since written a different paper, under the title ``Housing First or Treatment First? Ask the Kids, Not the Parents,'' that addresses both, and I would like to ask whether you would consider it as a new submission. The question in the new paper, the model, the experiments and the conclusions are all new. Among the changes, the model now includes children and the externalities a parent's drug use imposes on them. I use a simple, stripped-down two-period model to show precisely that uninsurable risk due to a substance use disorder (SUD) is the reason drug use is concentrated at the bottom of the income distribution, without having to assume, as the literature does, preferences that differ systematically with income.

\textbf{The advance.} The first paper asked whether conditioning housing assistance on sobriety reduces SUD and homelessness through moral hazard. It found that, with rationed vouchers, conditionality lowers SUD but raises homelessness and lowers welfare; Treatment First came out ahead only with universal vouchers and with a cost of premature death to survivors (loss of consortium) added to welfare as an exogenous parameter. The new paper drops that term and shows that moral hazard alone does not justify conditionality: without effects on children, a newborn is worse off under Treatment First, by 1.09 percent of lifetime consumption, even though drug use falls and the drug test's savings buy vouchers for other households. Instead, I show that the answer depends on the impact that Treatment First has on the children through its impact on drug use. It comes through three externalities. A parent's drug use costs the child consumption and housing during childhood, earnings as an adult, but a voucher revoked exposes the child to childhood homelessness. The parent only counts a part of these effects. Within the existing rationed programme, Treatment First cuts drug use by 21 percent and homelessness by 11 percent, and childhood homelessness, the externality its critics believe it creates, by 11 percent as well. It raises a newborn's welfare by 2.57 percent, three quarters of it an externality the parent does not count, and in general equilibrium the parents gain too, by 0.69 percent of their own consumption, because the savings from the reduction in SUD and voucher spending can be spent on more vouchers to sober households. The model reproduces many micro findings of randomised controlled trials, including that households with a disorder lose from the condition. In the model, about three quarters of them are worse off. What reverses the verdict are the macro effects that trials do not control for: specifically, the effects on the households the condition deters from using, the vouchers its savings fund, and the children.

\textbf{The empirical base.} Each mechanism that drives the answer is now pinned to direct evidence. The earnings penalty a child bears from a parent's SUD during their childhood is set to the natural-experiment estimate of Nilsson (2017), the deterrence of the drug test under Treatment First is set based on the US military's random testing (Mehay and Pacula, 1999), and the effect of housing on recovery to the randomised trial of Milby et al.\ (2005). Without being targeted, the model reproduces that trial's effect of making housing conditional on abstinence (1.40 against 1.26) and the effect of Housing First on the three-year mortality of homeless households with mental illness (veterans in the study) (3.6 against 4.6 percentage points; Costantini, 2025).

\textbf{The model and the transition.} The model also now includes other new elements: serious mental illness, a choice of treatment, a drug price that clears its market (it was held fixed in the first paper), and a budget balanced through the rationing of vouchers rather than through taxes. I also included transitions, which cover every cohort alive at the reform and clear the drug price and the budget every period. In the main transition, homelessness and childhood homelessness never rise, and children born at the reform gain 1.45 percent. They do not rise because few households lose their vouchers, due to low detection and households entering treatment, drug use falls from the first year, and the screen's savings fund about 7 percent more vouchers for others. Alternative transition specifications, in which households are unable to stop using, include an announced screen or a temporary exemption in treatment, and show that the impact of the reform on homelessness and childhood homelessness in the first years depends on the reform's design: catching every user with no exemption raises homelessness by up to 14 percent at first, while a five-period exemption for households in treatment holds the rise to 1.6 percent and keeps childhood homelessness below its Housing First level.

The question is current: in July 2025 an executive order directed federal agencies to end support for Housing First. Both papers are attached to this email. I would be grateful for your guidance.

Sincerely,\\
Shahar Rotberg\\
Canada Mortgage and Housing Corporation

\end{document}
